\begin{tikzpicture}[
    x=1pt,
    y=1pt,
    font=\small,
    node distance=16pt,
    >=stealth,
    stage/.style={
        draw=#1!75!black,
        fill=#1!10,
        rounded corners=10pt,
        line width=0.9pt,
        minimum width=106pt,
        minimum height=64pt,
        align=center
    },
    subtag/.style={
        draw=#1!70!black,
        fill=white,
        rounded corners=6pt,
        line width=0.8pt,
        inner xsep=7pt,
        inner ysep=4pt,
        font=\scriptsize
    },
    flow/.style={
        ->,
        line width=1pt,
        draw=gray!75
    },
    contextline/.style={
        ->,
        line width=0.9pt,
        draw=teal!70!black,
        dashed
    }
]
\def\W{730}
\def\H{248}

\fill[gray!4, rounded corners=16pt] (0,0) rectangle (\W,\H);
\draw[gray!35, rounded corners=16pt, line width=1pt] (0,0) rectangle (\W,\H);

\fill[cyan!18, rounded corners=14pt] (18,200) rectangle (712,230);
\node[anchor=west, font=\bfseries\small, text=black] at (34,215)
{Large Language Model Generation Pipeline};
\node[anchor=east, font=\scriptsize, text=gray!70!black] at (694,215)
{context-conditioned autoregressive decoding};

\node[stage=cyan] (prompt) at (78,118) {
    \textbf{输入提示}\\[4pt]
    用户问题\\
    系统指令\\
    历史消息
};

\node[stage=blue] (token) at (198,118) {
    \textbf{分词编码}\\[4pt]
    Text $\rightarrow$ Tokens\\
    位置编码\\
    向量表示
};

\node[stage=teal, minimum width=158pt, minimum height=84pt] (transformer) at (372,118) {
    \textbf{Transformer 堆叠}\\[4pt]
    多层自注意力\\
    前馈网络\\
    残差与归一化
};

\node[subtag=teal] (context) at (372,66) {上下文窗口};
\node[subtag=teal] (attention) at (372,169) {注意力计算};

\node[stage=orange] (prob) at (544,118) {
    \textbf{下一词概率分布}\\[4pt]
    logits\\
    Softmax\\
    候选词概率排序
};

\node[subtag=orange] (sampling) at (544,66) {概率采样};

\node[stage=red] (decode) at (650,118) {
    \textbf{自回归采样}\\[4pt]
    Top-k / Top-p\\
    温度控制\\
    生成下一个词元
};

\node[stage=violet] (output) at (650,118-92) {
    \textbf{输出回复}\\[4pt]
    逐词展开\\
    语义连贯\\
    最终响应
};

\draw[flow] (prompt.east) -- (token.west);
\draw[flow] (token.east) -- (transformer.west);
\draw[flow] (transformer.east) -- (prob.west);
\draw[flow] (prob.east) -- (decode.west);
\draw[flow] (decode.south) -- (output.north);

\draw[contextline] (prompt.north) .. controls (116,182) and (314,188) .. (transformer.north west);
\draw[contextline] (token.north) .. controls (228,180) and (326,180) .. (attention.west);
\draw[contextline] (context.east) .. controls (430,66) and (492,66) .. (sampling.west);

\draw[flow, rounded corners=10pt] (decode.west) |- (592,30) -| node[pos=0.55, below=8pt, font=\scriptsize, text=gray!70!black] {将新词元拼接回上下文并继续预测} (token.south);

\node[font=\scriptsize, align=left, text=gray!75!black] at (108,30) {
    核心机制: 在给定上下文 $x_{1:n}$ 条件下，模型持续预测 $P(x_{n+1}\mid x_{1:n})$
};
\end{tikzpicture}
